\chapter{1979 Nobel Prize in Economics}
\textbf{Laureates: }Theodore Schultz (USA), Arthur Lewis (UK)

\section*{Reason for Award}
For their pioneering research into economic development research with particular consideration of
the problems of developing countries

\section{What They Did}
Theodore Schultz emphasized the importance of human capital in agriculture and development in poor
countries, arguing that investments in education, health, and skills are productive forms of
investment alongside physical capital. He analyzed how the returns to these investments depend on
their costs, opportunities, and institutional setting.

Schultz's work on agricultural economics showed that traditional agriculture can be efficient given
the technologies, prices, and incentives available. He argued that modernization requires access
to new inputs, education, research, and institutional change. Schultz helped develop and apply the
human-capital approach, showing how education, training, and health can affect productivity and
economic welfare.

His analysis of rural development emphasized human capabilities and agricultural productivity.
Arthur Lewis developed the dual-sector model of economic development, explaining one possible
pattern of structural transformation from a traditional agricultural sector to a modern industrial
sector through the movement of labor. Lewis's model described how an economy with an elastic supply
of labor could industrialize before labor scarcity eventually increased wages.

He showed how a growing industrial sector could absorb surplus labor at relatively constant wages
until a turning point at which labor scarcity increased wages. Lewis also contributed to growth
theory, analyzing savings, investment, trade, and the constraints facing developing economies.
Lewis received the Nobel Prize for his work on development economics, emphasizing the importance
of structural transformation.

Schultz's work on human capital showed how investments in education, health, and training can
contribute to productivity and economic growth. His analysis of agricultural development
emphasized technology adoption and the incentives facing rural producers. Lewis's dual-sector
model provided a framework for studying industrialization, urbanization, and the movement of labor
between sectors.

\section{Impact on Economic Thinking}
Their work shifted attention beyond physical capital accumulation toward human capabilities and
structural change, helping shape modern development economics. Schultz's emphasis on human capital
informed research on agricultural investment, rural development, and industrialization. He argued
that education, health, and skills should be treated as economically important investments.

Lewis's analysis described how a growing industrial sector could absorb agricultural surplus labor,
creating a process of structural transformation. Their work drew attention to the interaction
between agriculture and industry in development. The human-capital concept became central to
growth theory and to explanations of differences in productivity and income.

The dual-sector model provided a framework for understanding urbanization, migration, and structural
transformation that characterizes development. Schultz's work on agricultural development
challenged the view that traditional agriculture was necessarily inefficient, showing that farmers
respond to incentives and available opportunities. Lewis's model provided a theoretical framework
for analyzing transitions from agrarian to industrial economies. Their work contributed to the
development of development economics as a distinct field.

Schultz's work on human capital influenced education policy and labor-market analysis. Lewis's
dual-sector model provided insights into structural transformation in developing economies. Their
combined contributions helped establish development economics as a major subfield.

Both thinkers emphasized that development depends on the interaction of people, productive
resources, institutions, and structural change, although they approached that problem from
different perspectives.

\section{Modern Perspectives Today}
Human capital theory has become central to education policy, labor economics, and growth theory.
Lewis's dual-sector model continues to inform structural-transformation analysis. Modern
development indicators also incorporate human-development and capability approaches, including
dimensions that extend beyond the original human-capital framework.

Dual-sector models continue to inform analysis of rural--urban migration and industrialization
sequencing. Human capital returns estimation guides education spending decisions in developing
countries. Understanding the complementary roles of agriculture and industry remains relevant
today, especially in Sub-Saharan Africa and South Asia undergoing similar transitions.

Schultz's work on agricultural development informs contemporary debates about food security and
rural transformation. Lewis's migration framework continues to inform analysis of urbanization in
developing economies. Modern development economics has expanded to incorporate institutions,
governance, international trade, and political economy.

The human-capital approach has influenced research on health, education, and poverty reduction.
Lewis's dual-sector model has been extended to incorporate informal sectors and urban unemployment,
addressing limitations of the original framework. Contemporary development indicators include
human-development metrics that capture education, health, and living standards. Their work
continues to inform debates about poverty reduction and economic transformation.

Contemporary research on human capital, structural transformation, and economic development
continues to build on the foundations laid by Schultz and Lewis, demonstrating the enduring
relevance of their contributions to development economics.

\subsection*{Key References}
\begin{itemize}
    \item Schultz, T. W. (1961). ``Investment in Human Capital.''
    \textit{American Economic Review}.
    \item Lewis, W. A. (1954). ``Economic Development with Unlimited Supplies of Labour.''
    \textit{The Manchester School}.
\end{itemize}
